% !PRINTED 496
\weberpage{496}

%<<CONT
den Fall er\"ortert, dass die Co\"efficienten einem bestimmten K\"orper
$\Omega$ angeh\"oren. Wir machen jetzt die Annahme, dass dieser K\"orper
$\Omega$ ein \gesperrt{algebraischer Zahlk\"orper} sei, und betrachten also
Ausdr\"ucke $\varphi(x, y, z, \ldots)$, die als eine Summe von Gliedern der
Form
\[
\alpha x^r y^s z^t \ldots
\]
dargestellt sind, worin die Exponenten $r, s, t, \ldots$ positive oder
wenigstens nicht negative ganze Zahlen sind, w\"ahrend die Coëfficienten
$\alpha$ Zahlen in $\Omega$ bedeuten. Einen solchen Ausdruck
\weq{1}{\varphi(x, y, z, \ldots) = \sum \alpha x^r y^s z^t \ldots}
nennen wir eine \gesperrt{ganze Function in $\Omega$}. Wir nehmen den
Ausdruck immer so geordnet und zusammengefasst an, dass dieselbe
Combination der Exponenten $r, s, t, \ldots$ nicht zweimal darin vorkommt,
und nennen zwei solche Ausdr\"ucke nur dann einander gleich, wenn sie
dieselben Producte $x^r y^s z^t \ldots$, mit denselben Coëfficienten
behaftet, enthalten. \gesperrt{Eine ganze Function wird dann und nur dann gleich Null
gesetzt, wenn alle ihre Coëfficienten Null sind.}

Mehrere ganze Functionen geben durch Addition, Subtraction und
Multiplication immer wieder ganze Functionen. Nach Bd.\,I, \S.\,143, I. kann
man f\"ur die Variablen $x, y, z, \ldots$ solche rationale Zahlwerthe
setzen, dass eine oder eine beliebige Anzahl von gegebenen von Null
verschiedenen ganzen Functionen in $\Omega$ nicht verschwindende Zahlwerthe
(in $\Omega$) erhalten. Daraus ergiebt sich, \gesperrt{dass ein Product mehrerer
ganzer Functionen nur dann verschwindet, wenn einer seiner Factoren
verschwindet.}

Die Summe $r + s + t + \cdots$ der Exponenten in einem Gliede des Ausdrucks
(1) heisst der \gesperrt{Grad dieses Gliedes}, und der gr\"osste Werth, den
der Grad eines Gliedes mit nicht verschwindendem Coëfficienten in $\varphi$
annimmt, heisst der \gesperrt{Grad der Function $\varphi$}.

Der Grad eines Productes aus zweien oder mehreren ganzen Functionen ist
gleich der Summe der Grade der einzelnen Factoren.

Denn fasst man in jedem der Factoren die Summe der Glieder h\"ochsten
Grades zu einer homogenen Function zusammen, so erh\"alt man die Glieder
h\"ochsten Grades des Productes, wenn man alle diese homogenen Functionen
mit einander multiplicirt.
